\documentclass[12pt,a4paper]{article}
\usepackage[UTF8]{ctex}
\usepackage{amsmath,amssymb,amsthm}
\usepackage{geometry}
\usepackage{bm}
\usepackage{hyperref}

\geometry{margin=2.5cm}

\title{力矩标记：一正一负为什么是空白？取消力后锥为什么变大？}
\author{现代机器人学：第12章 抓取与操作\\
详细解释}
\date{}

\begin{document}

\maketitle

\section{问题1：为什么一正一负就是空白？}

\subsection{力矩标记的基本规则}

对于平面中的点$p$，给定wrench锥$\text{pos}(\{F_1, F_2, \ldots, F_n\})$，标记规则为：

\begin{enumerate}
    \item 如果\textbf{每个}$F_i$对$p$产生\textbf{非负}力矩：标记为$'+'$
    \item 如果\textbf{每个}$F_i$对$p$产生\textbf{非正}力矩：标记为$'-'$
    \item 如果\textbf{每个}$F_i$对$p$产生\textbf{零}力矩：标记为$\pm$
    \item 如果\textbf{至少一个}$F_i$对$p$产生\textbf{正}力矩，且\textbf{至少一个}$F_j$对$p$产生\textbf{负}力矩：标记为\textbf{空白}
\end{enumerate}

\subsection{一正一负的情况}

假设对于点$p$：
\begin{itemize}
    \item $F_1$对$p$产生正力矩：$m_{1z}(p) > 0$
    \item $F_2$对$p$产生负力矩：$m_{2z}(p) < 0$
\end{itemize}

\subsection{正线性组合的力矩}

考虑所有正线性组合：
\begin{equation}
F = k_1 F_1 + k_2 F_2, \quad k_1, k_2 \geq 0
\end{equation}

对点$p$产生的总力矩为：
\begin{equation}
m_z(p) = k_1 m_{1z}(p) + k_2 m_{2z}(p) = k_1 \cdot (\text{正数}) + k_2 \cdot (\text{负数})
\end{equation}

\subsection{关键观察}

由于$k_1, k_2 \geq 0$，我们可以调整它们的比例：

\begin{itemize}
    \item 如果$k_1$很大，$k_2$很小：$m_z(p) > 0$（正力矩）
    \item 如果$k_1$很小，$k_2$很大：$m_z(p) < 0$（负力矩）
    \item 如果$k_1 / k_2 = -m_{2z}(p) / m_{1z}(p)$：$m_z(p) = 0$（零力矩）
\end{itemize}

\textbf{结论}：正线性组合可以产生\textbf{任意符号}的力矩！

\subsection{为什么不能标记为单一符号？}

\begin{itemize}
    \item \textbf{不能标记为$'+'$}：因为可以产生负力矩
    \item \textbf{不能标记为$'-'$}：因为可以产生正力矩
    \item \textbf{不能标记为$\pm$}：因为可以产生非零力矩
    \item \textbf{只能标记为空白}：表示力矩符号取决于系数选择，没有一致的符号
\end{itemize}

\subsection{标签的含义}

标签的含义是：\textbf{所有}正线性组合都产生\textbf{相同符号}的力矩。

对于一正一负的点：
\begin{itemize}
    \item 某些正线性组合产生正力矩
    \item 某些正线性组合产生负力矩
    \item 某些正线性组合产生零力矩
    \item \textbf{没有一致的符号}，因此不能分配单一标签
\end{itemize}

\section{问题2：为什么取消力后wrench锥变大了？}

\subsection{反直觉的现象}

\textbf{直觉}：减少力的数量应该减少可能的wrench范围。

\textbf{实际}：取消某些力后，wrench锥可能\textbf{变大}！

\subsection{关键理解：一致标记 vs 不一致标记}

\subsubsection{一致标记区域（有标签）}

\begin{itemize}
    \item 所有力对区域内点产生\textbf{相同符号}的力矩
    \item 正线性组合\textbf{只能}产生该符号的力矩
    \item 这是wrench锥的\textbf{限制}
    \item 标记为$'+'$、$'-'$或$\pm$
\end{itemize}

\subsubsection{不一致标记区域（空白）}

\begin{itemize}
    \item 不同力对区域内点产生\textbf{不同符号}的力矩
    \item 正线性组合可以产生\textbf{任意符号}的力矩
    \item 这是wrench锥的\textbf{灵活性}
    \item 标记为空白
\end{itemize}

\subsection{取消力的影响}

\subsubsection{情况1：取消产生一致标记的力}

假设有$F_1, F_2, F_3$三个力：

\begin{itemize}
    \item $F_1$和$F_2$对某个区域产生一致的正力矩（标记$'+'$）
    \item $F_3$对同一区域也产生正力矩
    \item 该区域保持标记$'+'$
\end{itemize}

如果取消$F_3$：
\begin{itemize}
    \item 该区域仍然由$F_1$和$F_2$产生一致标记
    \item 标记仍然为$'+'$
    \item \textbf{没有变化}
\end{itemize}

\subsubsection{情况2：取消产生不一致标记的力}

假设有$F_1, F_2, F_3$三个力：

\begin{itemize}
    \item $F_1$对某个区域产生正力矩
    \item $F_2$对同一区域产生负力矩
    \item $F_3$对同一区域产生正力矩
    \item 该区域标记为\textbf{空白}（不一致）
\end{itemize}

如果取消$F_3$：
\begin{itemize}
    \item 只剩下$F_1$（正）和$F_2$（负）
    \item 该区域仍然标记为\textbf{空白}（不一致）
    \item \textbf{没有变化}
\end{itemize}

\subsubsection{情况3：取消"限制性"的力}

这是关键情况！

假设有$F_1, F_2, F_3$三个力：

\begin{itemize}
    \item $F_1$和$F_2$对某个区域产生一致的正力矩（标记$'+'$）
    \item $F_3$对同一区域也产生正力矩
    \item 该区域标记为$'+'$（一致标记）
\end{itemize}

如果取消$F_3$：
\begin{itemize}
    \item 只剩下$F_1$和$F_2$
    \item 如果$F_1$和$F_2$对该区域产生\textbf{不一致}的力矩（一正一负）
    \item 该区域从$'+'$变为\textbf{空白}
    \item \textbf{限制减少了，灵活性增加了！}
\end{itemize}

\subsection{数学分析}

\subsubsection{Wrench锥的定义}

\begin{equation}
\mathcal{WC}_1 = \text{pos}(\{F_1, F_2, F_3\}) = \{k_1 F_1 + k_2 F_2 + k_3 F_3 \mid k_i \geq 0\}
\end{equation}

\begin{equation}
\mathcal{WC}_2 = \text{pos}(\{F_1, F_2\}) = \{k_1 F_1 + k_2 F_2 \mid k_i \geq 0\}
\end{equation}

\subsubsection{包含关系}

显然：
\begin{equation}
\mathcal{WC}_2 \subseteq \mathcal{WC}_1
\end{equation}

因为$\mathcal{WC}_2$中的任何wrench都可以通过设置$k_3 = 0$在$\mathcal{WC}_1$中表示。

\subsubsection{但是...}

\textbf{关键洞察}：虽然$\mathcal{WC}_2 \subseteq \mathcal{WC}_1$，但\textbf{灵活性}可能增加！

\subsubsection{灵活性的定义}

对于点$p$，考虑所有正线性组合产生的力矩范围：

\begin{itemize}
    \item \textbf{一致标记}：力矩范围受限（只能为正、只能为负、或只能为零）
    \item \textbf{不一致标记}：力矩范围不受限（可以为任意值）
\end{itemize}

\subsubsection{具体例子}

假设：
\begin{itemize}
    \item $F_1$对点$p$产生力矩：$m_{1z}(p) = 2$
    \item $F_2$对点$p$产生力矩：$m_{2z}(p) = -1$
    \item $F_3$对点$p$产生力矩：$m_{3z}(p) = 3$
\end{itemize}

\textbf{情况1}：使用$\{F_1, F_2, F_3\}$

总力矩：
\begin{equation}
m_z(p) = 2k_1 - k_2 + 3k_3
\end{equation}

由于$k_3 \geq 0$且$m_{3z}(p) = 3 > 0$：
\begin{itemize}
    \item 即使$k_1 = 0$，$k_2$很大，$k_3$也可以补偿
    \item 但所有三个力矩都是正或零的组合，可能产生一致标记
    \item 如果$F_1, F_2, F_3$对$p$都产生正力矩，标记为$'+'$
\end{itemize}

\textbf{情况2}：只使用$\{F_1, F_2\}$

总力矩：
\begin{equation}
m_z(p) = 2k_1 - k_2
\end{equation}

由于$k_1, k_2 \geq 0$：
\begin{itemize}
    \item 可以产生任意值的力矩（通过调整$k_1/k_2$）
    \item 标记为\textbf{空白}（不一致）
    \item \textbf{灵活性更大}！
\end{itemize}

\subsection{几何理解}

\subsubsection{Wrench空间中的表示}

在wrench空间中：

\begin{itemize}
    \item \textbf{一致标记区域}：对应wrench锥的\textbf{边界}或\textbf{限制方向}
    \item \textbf{不一致标记区域}：对应wrench锥的\textbf{内部}或\textbf{灵活方向}
\end{itemize}

\subsubsection{取消力的效果}

\begin{itemize}
    \item 取消某些力可能\textbf{移除限制}
    \item 使某些区域从一致标记变为不一致标记
    \item 虽然wrench锥的\textbf{集合}变小了（$\mathcal{WC}_2 \subseteq \mathcal{WC}_1$）
    \item 但在某些\textbf{方向}上的\textbf{灵活性}增加了
\end{itemize}

\section{实际例子}

\subsection{两个力的wrench锥}

考虑两个力$F_1$和$F_2$：

\begin{itemize}
    \item 如果它们对某个区域产生不一致的力矩
    \item 该区域标记为空白
    \item 可以产生任意符号的力矩
    \item \textbf{灵活性大}
\end{itemize}

\subsection{添加第三个力}

添加第三个力$F_3$：

\begin{itemize}
    \item 如果$F_3$对该区域也产生正力矩（与$F_1$和$F_2$中的正力矩一致）
    \item 该区域可能变为一致标记（$'+'$）
    \item 只能产生正力矩
    \item \textbf{灵活性减小}
\end{itemize}

\subsection{取消第三个力}

取消$F_3$：

\begin{itemize}
    \item 回到两个力的情况
    \item 如果$F_1$和$F_2$产生不一致的力矩
    \item 该区域变回空白
    \item \textbf{灵活性恢复}
\end{itemize}

\section{关键要点总结}

\subsection{问题1：一正一负为什么是空白？}

\begin{enumerate}
    \item 标签表示\textbf{所有}正线性组合都产生\textbf{相同符号}的力矩
    \item 一正一负意味着可以产生\textbf{任意符号}的力矩
    \item 没有一致的符号，因此不能分配单一标签
    \item 空白表示"不一致"，力矩符号取决于系数选择
\end{enumerate}

\subsection{问题2：取消力后锥为什么变大？}

\begin{enumerate}
    \item \textbf{集合大小}：wrench锥的集合确实变小了（$\mathcal{WC}_2 \subseteq \mathcal{WC}_1$）
    \item \textbf{灵活性}：但在某些方向上的灵活性可能增加
    \item \textbf{一致标记}：取消某些力可能移除限制，使区域从一致标记变为不一致标记
    \item \textbf{空白区域}：空白区域表示更大的灵活性（可以产生任意符号的力矩）
    \item \textbf{关键}：这是关于\textbf{灵活性}和\textbf{限制}的权衡，而不是简单的集合大小
\end{enumerate}

\section{物理直觉}

\subsection{约束与自由度的权衡}

\begin{itemize}
    \item \textbf{更多约束}：一致标记区域，力矩方向受限
    \item \textbf{更少约束}：不一致标记区域，力矩方向灵活
    \item \textbf{取消限制性约束}：可能增加灵活性
\end{itemize}

\subsection{力闭合的视角}

\begin{itemize}
    \item \textbf{力闭合要求}：没有一致标记区域（即没有限制）
    \item \textbf{更多力}：可能产生更多一致标记区域（限制）
    \item \textbf{更少力}：可能减少一致标记区域（增加灵活性）
    \item \textbf{但需要平衡}：足够的力来产生所需的wrench
\end{itemize}

\section{实际应用}

理解这些概念有助于：

\begin{itemize}
    \item \textbf{抓取设计}：理解如何平衡约束和灵活性
    \item \textbf{力闭合分析}：理解为什么某些配置可以实现力闭合
    \item \textbf{操作规划}：理解如何选择接触配置
    \item \textbf{优化设计}：在约束和灵活性之间找到平衡
\end{itemize}

\end{document}

